\section{Проблема манипуляции наградой и потенциал эволюционных алгоритмов в обеспечении надёжности ИИ-систем}

\subsection{Ключевые идеи}

Статья~\cite{asmussen2023reward} посвящена исследованию возможности использования эволюционных алгоритмов для решения проблемы манипуляции наградой.  
Манипуляция наградой возникает, когда агент находит способ увеличить свою награду, обходя изначально поставленную задачу.  
Такое поведение делает систему искусственного интеллекта непредсказуемой и небезопасной.

Основная цель работы~-- исследовать, могут ли эволюционные алгоритмы находить безопасные решения даже при наличии возможности манипуляции наградой.  

Рассматриваются два типа таких манипуляций: изменение функции награды и подделка входных данных.  

Авторы также сравнивают два подхода к описанию поведения агентов:
\begin{itemize}
    \item \emph{Общий подход}~--  не требует знаний о задаче, применим для любых сред.
    \item \emph{Специализированный подход (SPC)}~-- использует знания эксперта для выбора релевантных характеристик поведения.
\end{itemize}

Гипотеза авторов заключается в том, что сохранение разнообразия поведения повышает вероятность существования безопасных агентов даже при наличии манипулировании вознаграждения.  

\subsection{Математическое описание метода}

Каждый агент представляется нейронной сетью, параметры которой оптимизируются по функции приспособленности $R$.  
Вводится показатель безопасности $S$, характеризующий долю награды, полученной честным способом:
\begin{equation}
S \leq R
\end{equation}

Если $S = R$, поведение агента безопасно; если $S < R$, наблюдается манипуляция наградой.

Формы манипуляции:
\begin{itemize}
    \item изменение функции награды;
    \item изменение входных данных к функции награды.
\end{itemize}

Пример функции награды для среды «Камни и алмазы»:
\begin{equation}
R_t = r_D \cdot n_D - r_R \cdot n_R,
\end{equation}
где $n_D$ и $n_R$~-- количество алмазов и камней, $r_D$ и $r_R$~-- коэффициенты.  
Функция безопасности учитывает только корректные действия:
\begin{equation}
S_t = r_D^{(безопасн)} \cdot n_D^{(безопасн)} - r_R^{(безопасн)} \cdot n_R^{(безопасн)}.
\end{equation}

\subsection{Алгоритм}

Рассматривались четыре эволюционных алгоритма:

\begin{enumerate}
    \item \emph{NEAT}~-- классический алгоритм оптимизации структуры и весов нейросети.
    \item \emph{Поиск новизны (NS)}~-- максимизирует отличие поведения особей вместо награды.
    \item \emph{Поиск новизны с локальным соревнованием (NSLC)}~-- сочетает новизну и локальное сравнение по награде.
    \item \emph{Карта выдающихся особей (MAP-Elites)}~-- сохраняет лучших агентов в разных зонах поведенческого пространства.
\end{enumerate}

Новизна поведения для $i$-го агента вычисляется как среднее расстояние до $k$ ближайших соседей:
\begin{equation}
f_i = \frac{1}{k} \sum_{j=1}^{k} d(b_i, b_j),
\end{equation}
где $b_i$~-- вектор поведенческих характеристик, а $d(\cdot, \cdot)$~-- метрика расстояния.

\subsection{Наборы данных и среды}

Эксперименты проводились в двух средах из набора \emph{AI Safety Gridworlds}:

\begin{itemize}
    \item \emph{Камни и алмазы}~-- демонстрирует манипуляцию функцией награды.
    \item \emph{Полив томатов}~-- демонстрирует манипуляцию входными данными.
\end{itemize}

Среды представляют собой клеточные поля $10 \times 10$, где агент выполняет последовательность из 100 действий {вверх, вниз, влево, вправо}.

\subsection{Метрики}

Использовались две метрики:
\begin{itemize}
    \item \emph{Накопленная награда $R$ (fitness)}~-- суммарная награда, полученная агентом от среды в течение эпизода; отражает общую эффективность поведения, однако может расти и при нежелательном вмешательстве в вознаграждение, нарушая замысел задачи.
    
    \item \emph{Показатель безопасности $S$ (safety score)}~-- скрытая от алгоритма метрика, измеряющая долю награды, полученной корректным, безопасным способом. Если агент действует согласно целям разработчика, то $S = R$, а при манипуляции функцией вознаграждения или наблюдениями выполняется $S < R$.
\end{itemize}

\subsection{Модели и подходы}

Нейронные сети агентов принимали вектор наблюдений (56 признаков для «Камней и алмазов», 63~-- для «Полива томатов»).  
Выход сети определял действие:
\begin{equation}
a_t = \arg\max_i (y_i),
\end{equation}
где $y_i$~-- выходные значения сети.

Популяция содержала 200 агентов, каждый эксперимент повторялся 10 раз (всего $10^6$ итераций).  

Для каждого алгоритма применялись два варианта описания поведения:
\begin{enumerate}
    \item \emph{Специализированное}~-- характеристики поведения основаны на знании конкретной среды.
    \item \emph{Общее}~-- характеристики не зависят от задачи и формируются автоматически.
\end{enumerate}

\subsection{Результаты}

\begin{itemize}
    \item Классический NEAT не способен сохранять безопасные решения, так как максимизация награды приводит к манипуляции.
    \item Алгоритмы, поддерживающие поведенческое разнообразие (NS, NSLC, MAP-Elites), демонстрируют лучшие показатели безопасности.
    \item MAP-Elites обеспечивает наибольшее разнообразие стратегий и чаще сохраняет безопасных агентов.
    \item Для \emph{специализированного подхода} показатели безопасности выше, особенно в сложной среде «Полив томатов».
    \item \emph{Общий подход} даёт сопоставимые результаты в простой среде «Камни и алмазы», но значительно ухудшается при усложнении задачи.
\end{itemize}
